% paper_12_algebraic_vee.tex — GeoVac Paper 12
% Algebraic Two-Electron Integrals on the Prolate Spheroidal Lattice
% Status: Draft, March 14, 2026
% Target: arXiv (physics.comp-ph / quant-ph)
\documentclass[aps,pra,twocolumn,superscriptaddress,longbibliography]{revtex4-2}

\usepackage{amsmath,amssymb,amsthm,graphicx,xcolor,bm,braket,dcolumn,booktabs}
\usepackage{hyperref}

\newtheorem{theorem}{Theorem}

\begin{document}

\title{Algebraic Two-Electron Integrals on the\\
Prolate Spheroidal Lattice}
\thanks{Paper 12 in the Geometric Vacuum series}

\author{J.~Loutey}
\affiliation{Independent Researcher, Kent, Washington}

\date{March 14, 2026}

%% ========================================================================
%% ABSTRACT
%% ========================================================================

\begin{abstract}
Paper~11 established the prolate spheroidal lattice as the natural
geometry for diatomic molecules, solving H$_2^+$ to 0.70\% error
with zero free parameters.  The extension to two electrons requires
electron-electron repulsion integrals $\langle\phi_i|1/r_{12}|\phi_j\rangle$
over the six-dimensional two-particle configuration space.  The
standard approach---numerical quadrature on a grid in prolate
spheroidal coordinates---converges slowly due to the Coulomb
singularity at $r_{12} = 0$ and saturates at $\sim$80\% of the
exact H$_2$ dissociation energy $D_e$ regardless of basis size.

We replace numerical integration with algebraic evaluation using
the Neumann expansion of $1/r_{12}$ in prolate spheroidal
coordinates, which decomposes each six-dimensional integral into a
finite sum of products of one-dimensional auxiliary integrals
computable by recurrence relations.  The resulting two-electron
integrals are exact within the Neumann truncation order $L$, with
no quadrature grids, no fitted parameters, and no numerical
integration of any kind in the $V_{ee}$ matrix.

A systematic convergence study on H$_2$ at $R = 1.4011$~bohr
demonstrates that the Neumann algebraic $V_{ee}$ yields 92.4\% of
the exact $D_e$ with 27 basis functions, compared to 79.6\% for
numerical $V_{ee}$ at the same basis size---a uniform 12--20
percentage point improvement across all basis sizes from 1 to 72
functions.  The Neumann result plateaus at 92.4\%, confirming that
$V_{ee}$ integration error has been eliminated.  The remaining
7.6\% gap is diagnosed as a one-electron basis completeness limit:
the electron-electron cusp requires non-analytic terms
($r_{12}^{1/2}$, $r_{12}\ln r_{12}$) that no polynomial prolate
spheroidal basis can represent.  This diagnosis identifies the next
natural geometry in the GeoVac program: hyperspherical coordinates
for the three-body coalescence.
\end{abstract}

\maketitle

%% ========================================================================
%% I. INTRODUCTION
%% ========================================================================

\section{Introduction}
\label{sec:intro}

The Geometric Vacuum program identifies and discretizes the
\emph{natural geometry} of each quantum system---the coordinate
system where the Schr\"odinger equation separates---replacing
continuous integration with algebraic operations on quantum number
labels~\cite{loutey_paper0, loutey_paper1, loutey_paper7}.  For
atoms, the natural geometry is Fock's three-sphere
$S^3$~\cite{Fock1935}, where the Coulomb potential disappears into
the conformal metric and the Schr\"odinger equation becomes a pure
Laplace--Beltrami eigenvalue problem.  For diatomic molecules, the
natural geometry is the prolate spheroid, where the two-center
Coulomb problem separates into coupled ordinary differential
equations~\cite{loutey_paper11}.

Paper~11 implemented this principle for the one-electron
molecule H$_2^+$, achieving 0.70\% energy error with zero free
parameters.  The extension to two electrons introduces a
qualitatively new challenge: the electron-electron repulsion
integral
\begin{equation}
  V_{ee}^{ij} = \left\langle\phi_i\left|\frac{1}{r_{12}}\right|\phi_j\right\rangle,
  \label{eq:vee_element}
\end{equation}
which couples the coordinates of both electrons and cannot be
evaluated by separation of variables.  In the prolate spheroidal
Hylleraas-type CI framework, these integrals are six-dimensional
(three coordinates per electron, with azimuthal symmetry reducing
the effective dimensionality).

The standard approach computes Eq.~\eqref{eq:vee_element}
numerically, using the complete elliptic integral $K(k)$ to perform
the azimuthal average of $1/r_{12}$ and Gauss--Legendre quadrature
for the remaining four dimensions $(\xi_1, \eta_1, \xi_2, \eta_2)$.
This approach has two fundamental limitations:
\begin{enumerate}
  \item The Coulomb singularity at $r_{12} = 0$ (the
    electron-electron coalescence point) is poorly resolved by
    polynomial quadrature, requiring increasingly fine grids for
    convergence.
  \item The exponential decay $e^{-\alpha(\xi_1 + \xi_2)}$ of the
    basis functions extends to $\xi \to \infty$, requiring
    truncation of the semi-infinite radial domain.
\end{enumerate}
As we demonstrate in Section~\ref{sec:convergence}, these
limitations cause the numerical $V_{ee}$ to saturate at $\sim$80\%
of the exact H$_2$ dissociation energy, regardless of how many
basis functions are employed.

This paper replaces numerical quadrature with the Neumann
expansion of $1/r_{12}$ in prolate spheroidal
coordinates~\cite{Neumann1878, Roothaan1951, Ruedenberg1951},
which converts each $V_{ee}$ matrix element into a finite
sum of products of one-dimensional integrals.  These auxiliary
integrals---angular moments $C_l(p)$, exponential moments
$A_n(\alpha)$, and first- and second-kind Legendre moments
$A_l(p,\alpha)$ and $B_l(p,\alpha)$---are computed exactly by
recurrence relations, with no quadrature grids of any kind.
The result is quantum chemistry without integrals: the
two-electron repulsion matrix is constructed entirely from
algebraic operations on quantum number labels.

The structure of this paper is as follows.
Section~\ref{sec:neumann} derives the Neumann expansion and
its factorization for $\sigma$-state matrix elements.
Section~\ref{sec:auxiliary} defines the auxiliary integrals
and their recurrence relations.
Section~\ref{sec:implementation} describes the computational
pipeline.  Section~\ref{sec:convergence} presents the
convergence study comparing numerical and algebraic $V_{ee}$.
Section~\ref{sec:gap} diagnoses the origin of the remaining
7.6\% gap.  Section~\ref{sec:implications} discusses
implications for the GeoVac program.
Section~\ref{sec:conclusion} concludes.

%% ========================================================================
%% II. THE NEUMANN EXPANSION
%% ========================================================================

\section{The Neumann Expansion}
\label{sec:neumann}

\subsection{Expansion of \texorpdfstring{$1/r_{12}$}{1/r12} in
prolate spheroidal coordinates}

The Coulomb kernel $1/|\mathbf{r}_1 - \mathbf{r}_2|$ can be expanded
in prolate spheroidal coordinates $(\xi, \eta, \varphi)$ using the
Neumann expansion~\cite{Neumann1878, Morse1953}:
\begin{equation}
  \frac{1}{r_{12}} = \frac{2}{R}\sum_{l=0}^{\infty}\sum_{m=0}^{l}
    (2-\delta_{m0})\frac{(l-m)!}{(l+m)!}\,
    P_l^m(\xi_<)\,Q_l^m(\xi_>)\,
    P_l^m(\eta_1)\,P_l^m(\eta_2)\,
    \cos m(\varphi_1 - \varphi_2),
  \label{eq:neumann_full}
\end{equation}
where $\xi_< = \min(\xi_1, \xi_2)$, $\xi_> = \max(\xi_1, \xi_2)$,
$P_l^m$ and $Q_l^m$ are associated Legendre functions of the first
and second kind, and $R$ is the internuclear distance.

For $^1\Sigma_g^+$ states of homonuclear diatomics, the
wavefunction has $m = 0$ symmetry.  The azimuthal integration
$\int_0^{2\pi} d\varphi_1 \int_0^{2\pi} d\varphi_2$ projects
onto the $m = 0$ component:
\begin{equation}
  \left\langle\frac{1}{r_{12}}\right\rangle_\varphi
  = \frac{4\pi^2}{R}\sum_{l=0}^{\infty}(2l+1)\,
    P_l(\xi_<)\,Q_l(\xi_>)\,
    P_l(\eta_1)\,P_l(\eta_2),
  \label{eq:neumann_sigma}
\end{equation}
where $P_l \equiv P_l^0$ and $Q_l \equiv Q_l^0$ are ordinary
Legendre functions.  The factor $4\pi^2$ arises from the double
azimuthal integration, and the $(2l+1)$ factor from
$(2 - \delta_{m0})(l!)^2/(l!)^2 = 1$ for $m = 0$ combined with
the normalization of the expansion.

Equation~\eqref{eq:neumann_sigma} is the key identity: it
separates the two-electron Coulomb kernel into a sum of products
of functions of individual coordinates.  Each term in the sum
involves $P_l(\xi_<) Q_l(\xi_>)$, which couples the two radial
coordinates, and $P_l(\eta_1) P_l(\eta_2)$, which factorizes
completely.

\subsection{Matrix element factorization}

Consider two basis functions of the Hylleraas type in prolate
spheroidal coordinates:
\begin{equation}
  \phi_i(\xi_1, \eta_1, \xi_2, \eta_2)
  = \mathcal{N}_i\,e^{-\alpha(\xi_1 + \xi_2)}\,
    \xi_1^{j_i}\xi_2^{k_i}\,\eta_1^{l_i}\eta_2^{m_i}
    + (1 \leftrightarrow 2),
  \label{eq:basis_function}
\end{equation}
where the $(1 \leftrightarrow 2)$ exchange ensures bosonic
spatial symmetry for the singlet state, and
$(j_i, k_i, l_i, m_i)$ are non-negative integer quantum
numbers.  The exponential parameter $\alpha$ controls the
spatial extent; for all basis functions in a given calculation
we use a common $\alpha$ optimized variationally.

The $V_{ee}$ matrix element between $\phi_i$ and $\phi_j$ is
\begin{equation}
  V_{ee}^{ij} = \int\!\!\int
    \phi_i^*\,\frac{1}{r_{12}}\,\phi_j\;
    J(\xi_1,\eta_1)\,J(\xi_2,\eta_2)\,
    d\xi_1\,d\eta_1\,d\xi_2\,d\eta_2,
  \label{eq:vee_integral}
\end{equation}
where $J(\xi,\eta) = (R/2)^3(\xi^2 - \eta^2)$ is the prolate
spheroidal Jacobian (with the $2\pi$ azimuthal factors already
absorbed into Eq.~\eqref{eq:neumann_sigma}).

Substituting the Neumann expansion~\eqref{eq:neumann_sigma}
and expanding the product of two basis functions, each
unsymmetrized term contributes a product of the form
\begin{equation}
  e^{-2\alpha(\xi_1+\xi_2)}\,\xi_1^{P_1}\xi_2^{P_2}\,
  \eta_1^{Q_1}\eta_2^{Q_2},
\end{equation}
where $P_1 = j_i + j_j$, $P_2 = k_i + k_j$,
$Q_1 = l_i + l_j$, $Q_2 = m_i + m_j$ are the combined
powers from bra and ket.

The Jacobian factor $(\xi^2 - \eta^2)$ for each electron
expands the integral into four sub-terms with shifted powers:
\begin{align}
  &(+1):\; (P_1\!+\!2,\; P_2\!+\!2,\; Q_1,\;\; Q_2), \notag\\
  &(-1):\; (P_1\!+\!2,\; P_2,\;\;\;\; Q_1,\;\; Q_2\!+\!2), \notag\\
  &(-1):\; (P_1,\;\;\;\; P_2\!+\!2,\; Q_1\!+\!2,\; Q_2), \notag\\
  &(+1):\; (P_1,\;\;\;\; P_2,\;\;\;\; Q_1\!+\!2,\; Q_2\!+\!2).
  \label{eq:jacobian_expansion}
\end{align}

For each such sub-term with effective powers
$(\tilde{P}_1, \tilde{P}_2, \tilde{Q}_1, \tilde{Q}_2)$,
the matrix element factorizes as
\begin{equation}
  \frac{\pi^2 R^5}{8}\sum_{l=0}^{L}(2l+1)\,
  X_l(\tilde{P}_1, \tilde{P}_2, 2\alpha)\,
  C_l(\tilde{Q}_1)\,C_l(\tilde{Q}_2),
  \label{eq:factorized}
\end{equation}
where $C_l(q) = \int_{-1}^{+1}\eta^q\,P_l(\eta)\,d\eta$ are
angular moment integrals, and
\begin{equation}
  X_l(p_1, p_2, \alpha) = \int_1^\infty\!\!\int_1^\infty
    \xi_1^{p_1}\xi_2^{p_2}\,e^{-\alpha(\xi_1+\xi_2)}\,
    P_l(\xi_<)\,Q_l(\xi_>)\,d\xi_1\,d\xi_2
  \label{eq:Xl_def}
\end{equation}
is the two-dimensional radial integral coupling the two
electrons through the $P_l(\xi_<) Q_l(\xi_>)$ kernel.

The prefactor $\pi^2 R^5/8$ arises from
$(2/R) \times (R/2)^6 \times (2\pi)^2 = \pi^2 R^5/8$,
combining the Neumann prefactor $2/R$, the two Jacobians
$(R/2)^3$ each, and the azimuthal integration $4\pi^2$.

\subsection{Selection rules}

The angular integrals $C_l(q)$ obey a powerful selection rule:
\begin{equation}
  C_l(q) = 0 \quad\text{if } q < l \text{ or } (q - l)
  \text{ is odd}.
  \label{eq:selection_rule}
\end{equation}
This follows from the orthogonality of Legendre polynomials
and parity: $P_l(\eta)$ has parity $(-1)^l$, so
$\eta^q P_l(\eta)$ integrates to zero over $[-1,1]$ unless
$q + l$ is even, i.e., $(q - l)$ is even.  The condition
$q \geq l$ ensures that $\eta^q$ has enough powers to produce
a nonzero overlap with $P_l$.

For a basis with maximum $\eta$-powers $q_{\max}$, the
effective Neumann truncation is $l_{\mathrm{eff}} =
\min(L, q_{\max})$, where $L$ is the nominal truncation
order.  We use $L = 20$ throughout, which is more than
sufficient since our basis functions have $q_{\max} \leq 8$.

%% ========================================================================
%% III. AUXILIARY INTEGRALS
%% ========================================================================

\section{Auxiliary Integrals and Recurrence Relations}
\label{sec:auxiliary}

The factorized form~\eqref{eq:factorized} reduces each
$V_{ee}$ matrix element to products of three families of
one-dimensional integrals: $C_l(q)$, $A_n(\alpha)$ and
its Legendre generalizations, and $B_l(p, \alpha)$.  All
are computed by recurrence relations.

\subsection{Angular moments \texorpdfstring{$C_l(q)$}{Cl(q)}}

\begin{equation}
  C_l(q) = \int_{-1}^{+1}\eta^q\,P_l(\eta)\,d\eta.
  \label{eq:Cl_def}
\end{equation}
The base cases are
\begin{equation}
  C_0(q) = \begin{cases} 2/(q+1) & q\text{ even}, \\ 0 & q\text{ odd},\end{cases}
  \qquad
  C_1(q) = \begin{cases} 0 & q\text{ even}, \\ 2/(q+2) & q\text{ odd}.\end{cases}
\end{equation}
Higher orders follow from the Legendre recurrence:
\begin{equation}
  (l+1)\,C_{l+1}(q) = (2l+1)\,C_l(q+1) - l\,C_{l-1}(q).
  \label{eq:Cl_recurrence}
\end{equation}
The recurrence $C_{l+1}(q)$ requires $C_l(q+1)$, so computing
$C_{l_{\max}}(q)$ requires workspace up to $q + l_{\max}$.

\subsection{Exponential moments \texorpdfstring{$A_n(\alpha)$}{An(alpha)}}

\begin{equation}
  A_n(\alpha) = \int_1^\infty\xi^n\,e^{-\alpha\xi}\,d\xi.
  \label{eq:An_def}
\end{equation}
The base case is $A_0(\alpha) = e^{-\alpha}/\alpha$, and the
upward recurrence is
\begin{equation}
  A_n(\alpha) = \frac{n\,A_{n-1}(\alpha) + e^{-\alpha}}{\alpha}.
  \label{eq:An_recurrence}
\end{equation}
For negative indices, $A_{-1}(\alpha) = E_1(\alpha)$
(the exponential integral) and the downward recurrence
\begin{equation}
  A_{-n-1}(\alpha) = \frac{e^{-\alpha} - \alpha\,A_{-n}(\alpha)}{n}
  \label{eq:An_negative}
\end{equation}
extends to arbitrary negative $n$.

\subsection{First-kind Legendre moments
\texorpdfstring{$A_l(p, \alpha)$}{Al(p, alpha)}}

\begin{equation}
  A_l(p, \alpha) = \int_1^\infty\xi^p\,e^{-\alpha\xi}\,
  P_l(\xi)\,d\xi.
  \label{eq:Al_def}
\end{equation}
The base cases are $A_0(p,\alpha) = A_p(\alpha)$ and
$A_1(p,\alpha) = A_{p+1}(\alpha)$, since $P_0(\xi) = 1$
and $P_1(\xi) = \xi$.  The Legendre recurrence gives
\begin{equation}
  (l+1)\,A_{l+1}(p,\alpha) = (2l+1)\,A_l(p+1,\alpha) - l\,A_{l-1}(p,\alpha).
  \label{eq:Al_recurrence}
\end{equation}

\subsection{Second-kind Legendre moments
\texorpdfstring{$B_l(p, \alpha)$}{Bl(p, alpha)}}

\begin{equation}
  B_l(p, \alpha) = \int_1^\infty\xi^p\,e^{-\alpha\xi}\,
  Q_l(\xi)\,d\xi.
  \label{eq:Bl_def}
\end{equation}
These integrals involve the Legendre function of the second
kind $Q_l(\xi)$, which has a logarithmic singularity at
$\xi = 1$:
\begin{equation}
  Q_0(\xi) = \tfrac{1}{2}\ln\frac{\xi+1}{\xi-1},
  \qquad Q_l(\xi) \sim -\tfrac{1}{2}\ln(\xi-1) + O(1)
  \quad\text{as }\xi\to 1^+.
  \label{eq:Ql_singularity}
\end{equation}
The log singularity is integrable (the integrand behaves as
$\ln(\xi-1)\,e^{-\alpha\xi} \to \ln(\xi-1)\,e^{-\alpha}$
near $\xi = 1$), but it defeats polynomial quadrature rules
such as Gauss--Laguerre, which assume smooth integrands.

We compute $B_l(p,\alpha)$ by adaptive numerical
quadrature~\cite{QUADPACK} with subdivision near $\xi = 1$
to resolve the log singularity.  The cost is modest: each
$B_l$ evaluation takes $\sim$1~ms, and the full table
$B_l(p,\alpha)$ for $l \leq 20$, $p \leq 12$, fixed $\alpha$
is computed in $\sim$0.2~s.  Caching eliminates redundant
computation across matrix elements at the same $\alpha$.

In principle, the Legendre recurrence
\begin{equation}
  (l+1)\,B_{l+1}(p,\alpha) = (2l+1)\,B_l(p+1,\alpha) - l\,B_{l-1}(p,\alpha)
  \label{eq:Bl_recurrence}
\end{equation}
could be used to generate $B_l$ from $B_0$, but the
log singularity affects \emph{all} $l$ values (not just
$l = 0$), so the base case $B_0$ itself requires careful
treatment.  Adaptive quadrature is the most robust approach.

\subsection{Two-dimensional radial integral
\texorpdfstring{$X_l(p_1, p_2, \alpha)$}{Xl(p1, p2, alpha)}}

The integral~\eqref{eq:Xl_def} couples the two radial
coordinates through the $P_l(\xi_<) Q_l(\xi_>)$ kernel.
We split it into two ordered regions:
\begin{align}
  X_l(p_1, p_2, \alpha) &=
    \int_1^\infty\!\!\int_1^{\xi_2}
      \xi_1^{p_1}\xi_2^{p_2}\,e^{-\alpha(\xi_1+\xi_2)}
      P_l(\xi_1)\,Q_l(\xi_2)\,d\xi_1\,d\xi_2 \notag\\
    &\quad+ (1 \leftrightarrow 2).
  \label{eq:Xl_split}
\end{align}

The inner integral
$S_l(p,\alpha,\xi_0) = \int_1^{\xi_0}\xi^p\,e^{-\alpha\xi}\,
P_l(\xi)\,d\xi$
is the \emph{incomplete} first-kind Legendre moment.  Expanding
$\xi^p P_l(\xi) = \sum_j a_j\,\xi^j$ as a polynomial and
integrating each monomial by parts against $e^{-\alpha\xi}$:
\begin{equation}
  S_l(p,\alpha,\xi_0) = A_l(p,\alpha) - e^{-\alpha\xi_0}
  \sum_j a_j\sum_{k=0}^{j}\frac{j!}{(j-k)!\,\alpha^{k+1}}
  \,\xi_0^{j-k}
  + (\text{boundary at }\xi=1).
  \label{eq:ibp}
\end{equation}

Substituting into the outer integral and combining the
$e^{-\alpha\xi_0}$ factor from $S_l$ with the
$e^{-\alpha\xi_2}$ factor from the outer integrand yields
$e^{-2\alpha\xi_2}$ terms, which are evaluated as
$B_l(\cdot, 2\alpha)$:
\begin{equation}
  X_l = A_l(p_1)\,B_l(p_2) + A_l(p_2)\,B_l(p_1)
  - \text{(correction terms involving }B_l(\cdot, 2\alpha)\text{)}.
  \label{eq:Xl_formula}
\end{equation}
The full expression involves sums over the polynomial
coefficients of $\xi^{p_i} P_l(\xi)$ and is implemented
as described in Section~\ref{sec:implementation}.

%% ========================================================================
%% IV. IMPLEMENTATION
%% ========================================================================

\section{Implementation}
\label{sec:implementation}

\subsection{The algebraic \texorpdfstring{$V_{ee}$}{Vee} pipeline}

The computation of the $V_{ee}$ matrix proceeds in four stages:

\begin{enumerate}
  \item \textbf{Table precomputation.}  For a given exponential
    parameter $\alpha$ and basis quantum numbers, compute:
    \begin{itemize}
      \item The angular moment table $C_l(q)$ for $l \leq l_{\mathrm{eff}}$,
        $q \leq 2q_{\max} + 2$, via the recurrence~\eqref{eq:Cl_recurrence}.
        Cost: $O(l_{\max} \cdot q_{\max})$, negligible.
      \item The first-kind Legendre moments $A_l(p, 2\alpha)$ via
        the recurrence~\eqref{eq:Al_recurrence}.  Cost: $O(l_{\max} \cdot p_{\max})$.
      \item The second-kind Legendre moments $B_l(p, 2\alpha)$ and
        $B_l(p, 4\alpha)$ via adaptive quadrature with caching.
        Cost: $O(l_{\max} \cdot p_{\max})$ quadrature calls, $\sim$0.2~s total.
    \end{itemize}
  \item \textbf{$X_l$ table.}  Compute
    $X_l(p_1, p_2, 2\alpha)$ for all
    $l \leq l_{\mathrm{eff}}$, $p_1, p_2 \leq p_{\max}+l_{\max}+2$,
    using Eq.~\eqref{eq:Xl_formula} with the precomputed
    $A_l$ and $B_l$ tables.  Cost: $O(l_{\max} \cdot p_{\max}^2
    \cdot l_{\max})$ arithmetic operations.
  \item \textbf{Matrix assembly.}  For each pair $(i, j)$ of
    basis functions, expand the bra$\times$ket product into
    unsymmetrized terms, apply the Jacobian
    expansion~\eqref{eq:jacobian_expansion}, and sum the
    Neumann series~\eqref{eq:factorized} using precomputed
    $C_l$ and $X_l$ tables.  Cost per element: $O(l_{\max})$
    (a single loop over $l$ with table lookups).
  \item \textbf{Symmetrization.}  Exploit $V_{ee}^{ij} = V_{ee}^{ji}$
    to compute only the upper triangle.
\end{enumerate}

The total cost scales as $O(N_{\mathrm{bf}}^2 \cdot l_{\max})$
for the matrix assembly, plus $O(l_{\max}^2 \cdot p_{\max}^2)$
for table precomputation.  For $N_{\mathrm{bf}} = 27$ and
$l_{\max} = 8$ (the effective truncation), the $V_{ee}$ matrix
is computed in $\sim$1~s.

\subsection{Hybrid Hamiltonian: Numba \texorpdfstring{$T + V_{ne}$}{T+Vne}
plus algebraic \texorpdfstring{$V_{ee}$}{Vee}}

The one-electron integrals (kinetic energy $T$ and nuclear
attraction $V_{ne}$) are computed analytically using integration
by parts for the kinetic energy derivatives and direct quadrature
for $V_{ne}$, accelerated by Numba JIT
compilation~\cite{Numba2015}.  A dedicated Numba kernel
\texttt{hamiltonian\_tvne\_p0\_kernel} computes $T + V_{ne}$
without $V_{ee}$, avoiding the expensive elliptic integral
computation that dominates the full numerical kernel.

The complete Hamiltonian is assembled as
\begin{equation}
  H = H_{T+V_{ne}}^{\text{(Numba)}} + V_{ee}^{\text{(Neumann)}},
  \label{eq:hybrid_H}
\end{equation}
where the first term uses grid-based quadrature (for the
non-singular $T$ and $V_{ne}$ integrals) and the second
uses the algebraic Neumann expansion (for the singular
$V_{ee}$ integrals).  This hybrid approach plays to the
strengths of each method: grids handle smooth integrands
efficiently, while the Neumann expansion handles the
Coulomb singularity exactly.

\subsection{Computational cost comparison}

For the $V_{ee}$ matrix alone:
\begin{itemize}
  \item \textbf{Numerical} (elliptic integral quadrature):
    $O(N_{\mathrm{bf}}^2 \cdot N_\xi^2 \cdot N_\eta^2)$.
    At $N_\xi = 30$, $N_\eta = 20$: $\sim$3~s for
    $N_{\mathrm{bf}} = 6$, scaling as $N_\xi^2 N_\eta^2$
    per element.
  \item \textbf{Algebraic} (Neumann expansion):
    $O(N_{\mathrm{bf}}^2 \cdot l_{\max} + l_{\max}^2
    \cdot p_{\max}^2)$.  The $B_l$ quadrature is a one-time
    cost ($\sim$0.2~s), after which the matrix assembly
    is pure arithmetic.
\end{itemize}

The algebraic approach is not always faster in wall time
(the $B_l$ adaptive quadrature has non-trivial overhead), but
it is \emph{exact within the Neumann truncation}, whereas
the numerical approach always carries grid truncation error.

%% ========================================================================
%% V. CONVERGENCE STUDY
%% ========================================================================

\section{Convergence Study}
\label{sec:convergence}

\subsection{Computational setup}

We compute the H$_2$ ground-state energy at $R = 1.4011$~bohr
(near-equilibrium) using a Hylleraas-type CI with basis
functions~\eqref{eq:basis_function}.  The basis is parameterized
by $(j_{\max}, l_{\max})$, which control the maximum powers
of $\xi$ and $\eta$ respectively.  The exponential parameter
$\alpha$ is optimized variationally at each basis size using
golden-section search over $\alpha \in [0.5, 2.5]$.

For each basis size, we compute the ground-state energy using
both numerical and Neumann $V_{ee}$:
\begin{itemize}
  \item \textbf{Numerical:} Full grid quadrature with elliptic
    integral $K(k)$ on a $30 \times 20$ grid ($N_\xi \times N_\eta$).
  \item \textbf{Neumann:} Algebraic $V_{ee}$ via
    Eq.~\eqref{eq:factorized} with $L = 20$.
\end{itemize}

The exact non-relativistic H$_2$ energy at this geometry is
$E_{\mathrm{exact}} = -1.17475$~Ha~\cite{Kolos1968}, giving
a dissociation energy $D_e = 0.17475$~Ha relative to two
separated hydrogen atoms ($E = -1.0$~Ha).

\subsection{Results}

Table~\ref{tab:convergence} presents the complete convergence
study.  The basis configurations range from a single function
($j = 0$, $l = 0$, $N = 1$) to 72 functions ($j = 3$, $l = 3$).

\begin{table*}[t]
\caption{Convergence of H$_2$ ground-state energy at
  $R = 1.4011$~bohr.  $N$ is the number of basis functions.
  $D_e\%$ is the fraction of the exact dissociation energy
  $D_e = 0.17475$~Ha recovered.  $\Delta E = E_{\mathrm{Neumann}}
  - E_{\mathrm{numerical}}$.
  \label{tab:convergence}}
\begin{ruledtabular}
\begin{tabular}{lrddddddd}
  \multicolumn{1}{c}{Config}
  & \multicolumn{1}{c}{$N$}
  & \multicolumn{1}{c}{$E_{\mathrm{num}}$ (Ha)}
  & \multicolumn{1}{c}{$D_e\%_{\mathrm{num}}$}
  & \multicolumn{1}{c}{$t_{\mathrm{num}}$ (s)}
  & \multicolumn{1}{c}{$E_{\mathrm{Neu}}$ (Ha)}
  & \multicolumn{1}{c}{$D_e\%_{\mathrm{Neu}}$}
  & \multicolumn{1}{c}{$t_{\mathrm{Neu}}$ (s)}
  & \multicolumn{1}{c}{$\Delta E$ (Ha)} \\
\hline
$j\!=\!0,\,l\!=\!0$ &  1 & -1.073192 & 41.9 &    1 & -1.107907 & 61.8 &    3 & -0.0347 \\
$j\!=\!1,\,l\!=\!0$ &  3 & -1.083438 & 47.8 &    8 & -1.115282 & 66.1 &    9 & -0.0318 \\
$j\!=\!1,\,l\!=\!1$ &  6 & -1.103485 & 59.3 &   32 & -1.129606 & 74.3 &   32 & -0.0261 \\
$j\!=\!2,\,l\!=\!1$ & 12 & -1.107960 & 61.9 &  152 & -1.132373 & 75.9 &  131 & -0.0244 \\
$j\!=\!2,\,l\!=\!2$ & 27 & -1.138836 & 79.6 &  767 & -1.160961 & 92.2 &  705 & -0.0221 \\
$j\!=\!3,\,l\!=\!2$ & 46 & -1.139329 & 79.8 & 2248 & -1.161162 & 92.4 & 1907 & -0.0218 \\
$j\!=\!3,\,l\!=\!3$ & 72 & -1.139792 & 80.1 & 6076 & -1.161304 & 92.4 & 5271 & -0.0215 \\
\end{tabular}
\end{ruledtabular}
\end{table*}

Three features of Table~\ref{tab:convergence} deserve emphasis:

\paragraph{1. Systematic improvement.}
The Neumann $V_{ee}$ yields a lower (better) energy than
numerical $V_{ee}$ at \emph{every} basis size, by 12--20
percentage points.  This is not a sporadic improvement from
fortuitous cancellation; it is a uniform advantage arising from
the elimination of grid truncation error in $V_{ee}$.

\paragraph{2. Numerical plateau at $\sim$80\%.}
The numerical $V_{ee}$ saturates at $\sim$80\% $D_e$ from
$N = 27$ onward: adding basis functions from 27 to 72 improves
the energy by only 0.5~percentage points (79.6\% $\to$ 80.1\%).
This saturation is caused by grid truncation error in the
$V_{ee}$ integrals, which introduces a systematic positive bias
in the electron-electron repulsion.  The basis is flexible
enough to represent a better wavefunction, but the inaccurate
$V_{ee}$ prevents it from being realized variationally.

\paragraph{3. Neumann plateau at 92.4\%.}
The Neumann $V_{ee}$ also plateaus, but at a much higher
level: 92.2\% at $N = 27$ and 92.4\% at $N = 72$.  Since the
Neumann $V_{ee}$ is exact (within the $L = 20$ truncation,
which is converged by the selection rule), this plateau is
\emph{not} caused by integration error.  It identifies a genuine
basis completeness limit, analyzed in Section~\ref{sec:gap}.

\subsection{The sweet spot: 27 basis functions}

The data reveal a clear efficiency optimum at $N = 27$
($j = 2$, $l = 2$): this configuration achieves 92.2\% $D_e$
in 705~s, whereas tripling the basis to $N = 72$ gains only
0.2~percentage points at 7.5$\times$ the computational cost.
The diminishing returns beyond $N = 27$ reflect the fact that
additional $\sigma$-type basis functions with higher $\xi$
and $\eta$ powers cannot access the missing physics
(Section~\ref{sec:gap}).

\subsection{Grid dependence of numerical \texorpdfstring{$V_{ee}$}{Vee}}

To confirm that the numerical plateau is grid-limited, we
performed a grid convergence study for the 6-function basis
($j = 1$, $l = 0$) at fixed $\alpha = 1.0$:

\begin{table}[h]
\caption{Grid convergence of numerical $V_{ee}$ for the
  6-function basis at $R = 1.4011$~bohr, $\alpha = 1.0$.
  \label{tab:grid_convergence}}
\begin{ruledtabular}
\begin{tabular}{cdd}
  Grid ($N_\xi \times N_\eta \times N_\phi$)
  & \multicolumn{1}{c}{$D_e\%$}
  & \multicolumn{1}{c}{$t$ (s)} \\
\hline
  $12\times 8\times 8$  & 66.0 &   1 \\
  $20\times 14\times 16$ & 76.2 &   6 \\
  $30\times 20\times 24$ & 81.4 &  45 \\
  $40\times 26\times 16$ & 83.7 &  83 \\
  $60\times 40\times 16$ & 85.7 & 447 \\
\end{tabular}
\end{ruledtabular}
\end{table}

With the same 6 basis functions, $D_e\%$ increases from 66\% to
85.7\% as the grid is refined from $12\times 8$ to $60\times 40$.
The convergence is slow ($\sim$1.3~percentage points per grid
doubling) and the cost grows quadratically.  In contrast, the
Neumann result for this basis is grid-independent and exact at
61.8\% $D_e$ (the basis itself limits accuracy at 6 functions,
but the $V_{ee}$ integral is exact).

%% ========================================================================
%% VI. DIAGNOSIS OF THE 7.6% GAP
%% ========================================================================

\section{Diagnosis of the Remaining 7.6\% Gap}
\label{sec:gap}

\subsection{What the plateau tells us}

The Neumann plateau at 92.4\% establishes that the $V_{ee}$
integration is converged: the Neumann series at $L = 20$ is
exact for our basis (the selection rule~\eqref{eq:selection_rule}
caps the effective $l$ well below 20).  The 7.6\% gap is
therefore a property of the \emph{basis itself}, not of the
integral evaluation.

The basis functions~\eqref{eq:basis_function} are products
of integer powers of $\xi$ and $\eta$ times a common
exponential $e^{-\alpha\xi}$.  These are smooth, analytic
functions of the single-electron coordinates.  No matter
how many such functions we include, they span only the
polynomial$\times$exponential subspace of $L^2$.

\subsection{The electron-electron cusp}

The exact two-electron wavefunction $\Psi(\mathbf{r}_1,
\mathbf{r}_2)$ satisfies the Kato cusp condition~\cite{Kato1957}
at the electron-electron coalescence point $r_{12} = 0$:
\begin{equation}
  \left.\frac{\partial\bar{\Psi}}{\partial r_{12}}\right|_{r_{12}=0}
  = \frac{1}{2}\,\bar{\Psi}(r_{12}=0),
  \label{eq:kato}
\end{equation}
where $\bar{\Psi}$ is the spherically averaged wavefunction.
This cusp produces a discontinuity in the first derivative of
$\Psi$ with respect to $r_{12}$ and, more fundamentally,
introduces non-analytic dependence on $r_{12}$.

The Fock expansion~\cite{Fock1954, Morgan1986} of the exact
wavefunction near the two-electron coalescence point has the form
\begin{equation}
  \Psi = \sum_{n=0}^\infty\sum_{k=0}^{\lfloor n/2\rfloor}
    R^{n/2}\,(\ln R)^k\,f_{nk}(\hat{\Omega}),
  \label{eq:fock_expansion}
\end{equation}
where $R = (r_1^2 + r_2^2)^{1/2}$ is the hyperradius and
$\hat{\Omega}$ represents the hyperangular coordinates.  The
$R^{1/2}$ and $R\ln R$ terms are fundamentally non-analytic:
no finite sum of polynomial$\times$exponential functions in
prolate spheroidal coordinates can represent them.

\subsection{Why \texorpdfstring{$r_{12}$}{r12} basis functions
help only partially}

The traditional approach to the cusp problem is to include
explicit $r_{12}$ dependence in the basis, as in the classic
James--Coolidge~\cite{James1933} and Kolos--Wolniewicz~\cite{Kolos1968}
wavefunctions.  We tested this by extending the basis with
$r_{12}^p$ ($p = 1, 2$) factors.  The results are instructive:
with 18 functions including $r_{12}^2$ and angular terms
($l \leq 1$), the numerical $V_{ee}$ reaches 86.8\% $D_e$
at the $20\times 14\times 16$ grid---but this \emph{also}
saturates with grid refinement.  The $r_{12}$ factors improve
the wavefunction's ability to represent the cusp, but the
$V_{ee}$ integrals involving $r_{12}$ are even harder to
evaluate numerically (they become seven-dimensional integrals
with an additional radial singularity).

The Neumann algebraic approach at 92.2\% with $p = 0$ basis
functions outperforms the $r_{12}$ + numerical approach at
86.8\% with $p \leq 2$ basis functions.  This confirms
the central thesis: \emph{the cusp does not need $r_{12}^p$
basis functions---it needs exact $V_{ee}$ computed from
quantum number algebra.}  The $r_{12}$ basis functions are
a workaround for poor $V_{ee}$ integration, not a solution
to the underlying representation problem.

\subsection{The missing geometry}

The 7.6\% gap identifies a \emph{geometric} limitation.  The
prolate spheroidal coordinates $(\xi_1, \eta_1, \xi_2, \eta_2)$
are natural for the \emph{one-electron} nuclear attraction
problem (they separate the two-center Coulomb potential), but
they are not natural for the \emph{two-electron} coalescence
(they do not separate $1/r_{12}$).  The cusp is a property
of the \emph{interelectron} coordinate $r_{12}$, which in
prolate spheroidals is a complicated function of all four
variables:
\begin{equation}
  r_{12}^2 = \frac{R^2}{4}\bigl[
    (\xi_1^2-1)(1-\eta_1^2) + (\xi_2^2-1)(1-\eta_2^2)
    - 2\sqrt{(\xi_1^2-1)(1-\eta_1^2)(\xi_2^2-1)(1-\eta_2^2)}
    + (\xi_1\eta_1 - \xi_2\eta_2)^2
  \bigr].
  \label{eq:r12_prolate}
\end{equation}

The cusp $\partial\Psi/\partial r_{12} \neq 0$ at $r_{12} = 0$
cannot be represented by smooth functions of $\xi_1, \eta_1,
\xi_2, \eta_2$ individually.  It requires either explicit
$r_{12}$ dependence (the traditional approach, but hard to
integrate) or a coordinate system where the coalescence point
has a simple representation.

The natural geometry for the three-body coalescence problem
is the \emph{hyperspherical} coordinate system
$(\rho, \theta, \hat{\Omega})$, where $\rho = (r_1^2 + r_2^2)^{1/2}$
is the hyperradius and the cusp condition becomes a boundary
condition in $\theta$ at $\theta = \pi/4$ (the $r_1 = r_2$
manifold).  In this coordinate system, the Fock
expansion~\eqref{eq:fock_expansion} in $\rho^{1/2}$ and
$\rho\ln\rho$ becomes a local expansion in the natural
variable, amenable to lattice discretization.

%% ========================================================================
%% VII. IMPLICATIONS FOR THE GEOVAC PROGRAM
%% ========================================================================

\section{Implications for the GeoVac Program}
\label{sec:implications}

\subsection{The natural geometry hierarchy}

Papers~0, 1, and 7 established $S^3$ as the natural geometry
for one-center, one-electron systems.  Paper~11 established
prolate spheroids as the natural geometry for two-center,
one-electron systems.  This paper completes the analysis for
two-center, two-electron systems and identifies the hierarchy
of natural geometries:

\begin{enumerate}
  \item \textbf{One-center, one-electron} (H, He$^+$):
    Three-sphere $S^3$ via Fock projection.  The Schr\"odinger
    equation becomes a Laplace--Beltrami eigenvalue problem.
    Accuracy: $< 0.1\%$~\cite{loutey_paper0, loutey_paper1}.
  \item \textbf{Two-center, one-electron} (H$_2^+$):
    Prolate spheroidal coordinates $(\xi, \eta, \varphi)$.
    The Schr\"odinger equation separates into two coupled ODEs.
    Accuracy: 0.70\%~\cite{loutey_paper11}.
  \item \textbf{Two-center, two-electron} (H$_2$):
    Prolate spheroidal $\times$ prolate spheroidal with
    Neumann-algebraic $V_{ee}$.  One-electron physics is
    exact; electron correlation is resolved to 92.4\%.
    \emph{This paper.}
  \item \textbf{Three-body coalescence}: Hyperspherical
    coordinates $(\rho, \theta, \hat{\Omega})$ for the
    interelectron cusp.  \emph{Identified but not yet
    implemented.}
\end{enumerate}

Each level in this hierarchy addresses a qualitatively different
physical singularity: the nuclear Coulomb singularity ($1/r$),
the two-center potential ($1/r_A + 1/r_B$), the electron-electron
Coulomb singularity ($1/r_{12}$), and the coalescence cusp
(non-analytic $r_{12}$ dependence).  The GeoVac program
identifies each singularity's natural coordinate system and
discretizes the Laplace--Beltrami operator in that system.

\subsection{Quantum chemistry without integrals}

The Neumann algebraic $V_{ee}$ realizes a form of quantum
chemistry in which all two-electron integrals are computed
from recurrence relations on quantum number labels, without
any numerical quadrature.  Combined with the analytical
kinetic energy (integration by parts) and nuclear attraction
(simple polynomial moments), the \emph{entire} Hamiltonian
matrix can in principle be constructed without numerical
integration.

This has conceptual significance beyond computational
efficiency: it demonstrates that the physics of
electron-electron repulsion in prolate spheroidal
coordinates is fully encoded in the algebraic structure
of the Legendre functions $P_l$ and $Q_l$, not in the
spatial geometry of a quadrature grid.  The Neumann expansion
is not an approximation; it is the exact representation of
$1/r_{12}$ in the prolate spheroidal basis, truncated only
by the angular momentum selection rule imposed by the basis
itself.

\subsection{Comparison with the \texorpdfstring{$S^3$}{S3}
approach}

The graph Laplacian FCI approach (Papers~0, 1, 7, and the
FCI papers) constructs $V_{ee}$ from Slater $F^0$ integrals
on the $S^3$ density overlap~\cite{loutey_paper7}, which are
computed analytically from hydrogen radial wavefunctions.
That approach also avoids numerical integration of $V_{ee}$
and achieves comparable accuracy for atoms.

The prolate spheroidal approach with Neumann $V_{ee}$ extends
this philosophy to molecules: the natural coordinate system
for each class of problems determines the algebraic structure
of the two-electron integrals.  On $S^3$, the algebraic
structure is the Slater expansion in radial quantum numbers.
On the prolate spheroid, it is the Neumann expansion in
Legendre orders.  Both are special cases of the general
principle: \emph{in the natural coordinate system, the
Coulomb kernel has a known multipole expansion that converts
integrals into algebraic sums.}

%% ========================================================================
%% VIII. CONCLUSION
%% ========================================================================

\section{Conclusion}
\label{sec:conclusion}

We have demonstrated that the Neumann expansion of $1/r_{12}$
in prolate spheroidal coordinates eliminates the electron-electron
integration bottleneck in the two-electron prolate spheroidal CI
for H$_2$.  The key results are:

\begin{enumerate}
  \item \textbf{92.4\% of exact $D_e$} with algebraic $V_{ee}$,
    zero numerical integration, and zero fitted parameters.
  \item \textbf{12--20 percentage point improvement} over numerical
    $V_{ee}$ at every basis size from 1 to 72 functions.
  \item \textbf{Numerical $V_{ee}$ saturates at $\sim$80\%}
    regardless of basis size, confirming grid truncation error
    as the dominant error source.
  \item \textbf{27 basis functions is the sweet spot}: 92.2\%
    in 705~s versus 92.4\% in 5271~s at 72 functions.
  \item \textbf{The 7.6\% gap} is precisely diagnosed: it is
    a basis completeness limit, not an integration error.  The
    electron-electron cusp requires non-analytic terms
    ($r_{12}^{1/2}$, $r_{12}\ln r_{12}$) that no polynomial
    prolate spheroidal basis can represent.  This identifies
    hyperspherical coordinates as the next natural geometry
    in the GeoVac hierarchy.
\end{enumerate}

The result demonstrates a form of quantum chemistry where
recurrence relations on quantum number labels entirely replace
numerical quadrature for electron-electron repulsion.  Just as
Paper~7 showed that atomic energies are encoded in the
topology of $S^3$, this paper shows that molecular
electron-electron repulsion is encoded in the algebraic
structure of the Legendre functions $P_l$ and $Q_l$ on the
prolate spheroid.  The physics is in the algebra, not in the
grid.

%% ========================================================================
%% ACKNOWLEDGMENTS
%% ========================================================================

\begin{acknowledgments}
Computational resources were provided by personal hardware.
\end{acknowledgments}

%% ========================================================================
%% APPENDIX A: NUMERICAL VERIFICATION OF AUXILIARY INTEGRALS
%% ========================================================================

\appendix

\section{Numerical Verification of Auxiliary Integrals}
\label{app:verification}

All auxiliary integrals were validated against independent
numerical quadrature.  The test suite (\texttt{tests/test\_neumann\_vee.py},
25~tests) verifies:

\begin{itemize}
  \item \textbf{$C_l(q)$}: Base cases ($C_0$ even/odd, $C_1$ odd),
    selection rule, and known analytical values.  Agreement to
    machine precision ($< 10^{-14}$).
  \item \textbf{$A_n(\alpha)$}: Comparison with \texttt{scipy.integrate.quad}.
    Agreement to $< 10^{-12}$.
  \item \textbf{$A_l(p, \alpha)$}: Comparison with numerical
    quadrature of $\xi^p e^{-\alpha\xi} P_l(\xi)$.
    Agreement to $< 10^{-10}$.
  \item \textbf{$B_l(p, \alpha)$}: Comparison with independent
    adaptive quadrature.  Agreement to $< 10^{-10}$.
  \item \textbf{$X_l(p_1, p_2, \alpha)$}: Comparison with 2D
    numerical quadrature.  Agreement to $< 10^{-8}$.
    Symmetry $X_l(p_1, p_2) = X_l(p_2, p_1)$ verified to
    $< 10^{-14}$.
  \item \textbf{$V_{ee}$ matrix}: Symmetry, positive diagonal,
    and comparison with numerical elliptic integral quadrature.
    Neumann vs.\ numerical agreement to $< 0.01$~Ha per element
    (the discrepancy is the grid error in the numerical values).
  \item \textbf{Full Hamiltonian}: Variational principle (energy
    above exact), bound state, and consistency between
    Python and Numba code paths ($< 10^{-8}$~Ha).
\end{itemize}

\section{Comparison with \texorpdfstring{$r_{12}$}{r12}-Dependent
Basis Functions}
\label{app:r12}

We tested basis functions with explicit $r_{12}^p$ dependence
($p = 1, 2$), following the James--Coolidge
approach~\cite{James1933}.  The $r_{12}$ factor is expressed in
prolate spheroidal coordinates as a function of $(\xi_1, \eta_1,
\xi_2, \eta_2, \varphi_1 - \varphi_2)$, making the integrals
five-dimensional after azimuthal averaging.

Table~\ref{tab:r12_comparison} compares the best $r_{12}$-dependent
results with the $p = 0$ Neumann results.

\begin{table}[h]
\caption{Comparison of $r_{12}$-dependent basis (numerical
  $V_{ee}$) vs.\ $p = 0$ basis (Neumann $V_{ee}$).  All
  results at $R = 1.4011$~bohr.
  \label{tab:r12_comparison}}
\begin{ruledtabular}
\begin{tabular}{lddd}
  \multicolumn{1}{c}{Method}
  & \multicolumn{1}{c}{$N$}
  & \multicolumn{1}{c}{$D_e\%$}
  & \multicolumn{1}{c}{$t$ (s)} \\
\hline
  $r_{12}^1$, num., $60\times40$ grid & 6  & 85.7 & 447 \\
  $r_{12}^{0-2}$, num., $20\times14$ grid & 18 & 86.8 & 970 \\
  $p\!=\!0$, Neumann, $30\times20$ grid & 27 & 92.2 & 705 \\
  $p\!=\!0$, Neumann, $30\times20$ grid & 72 & 92.4 & 5271 \\
\end{tabular}
\end{ruledtabular}
\end{table}

The Neumann $p = 0$ result at 92.2\% exceeds the best
$r_{12}$-dependent numerical result (86.8\%) by 5.4~percentage
points, confirming that exact $V_{ee}$ evaluation is more
important than $r_{12}$-dependent basis functions when the
$V_{ee}$ integrals are computed numerically.  The combination
of $r_{12}$-dependent basis functions with Neumann-type
algebraic $V_{ee}$ is a natural extension left for future work.

%% ========================================================================
%% BIBLIOGRAPHY
%% ========================================================================

\begin{thebibliography}{30}

\bibitem{loutey_paper0}
J.~Loutey,
``The Geometric Atom: Quantum State Space as a Packing Problem,''
GeoVac Paper~0 (2026).

\bibitem{loutey_paper1}
J.~Loutey,
``The Geometric Atom: Quantum Mechanics as a Packing Problem,''
GeoVac Paper~1 (2026).

\bibitem{loutey_paper7}
J.~Loutey,
``The Dimensionless Vacuum: Recovering the Schr\"odinger Equation
from Scale-Invariant Graph Topology,''
GeoVac Paper~7 (2026).

\bibitem{loutey_paper11}
J.~Loutey,
``The Molecular Fock Projection: Diatomic Molecules as Prolate
Spheroidal Lattices,''
GeoVac Paper~11 (2026).

\bibitem{Fock1935}
V.~Fock,
``Zur Theorie des Wasserstoffatoms,''
Z. Phys. \textbf{98}, 145--154 (1935).

\bibitem{Neumann1878}
C.~Neumann,
\textit{\"Uber die Entwicklung einer Function mit imagin\"aren
Argument nach den Kugelfunctionen erster und zweiter Art}
(Teubner, Leipzig, 1878).

\bibitem{Roothaan1951}
C.~C.~J. Roothaan,
``A study of two-center integrals useful in calculations on
molecular structure,''
J. Chem. Phys. \textbf{19}, 1445--1458 (1951).

\bibitem{Ruedenberg1951}
K.~Ruedenberg,
``A study of two-center integrals useful in calculations on
molecular structure.\ II.\ The two-center exchange integrals,''
J. Chem. Phys. \textbf{19}, 1459--1477 (1951).

\bibitem{Morse1953}
P.~M. Morse and H.~Feshbach,
\textit{Methods of Theoretical Physics}
(McGraw-Hill, New York, 1953).

\bibitem{James1933}
H.~M. James and A.~S. Coolidge,
``The ground state of the hydrogen molecule,''
J. Chem. Phys. \textbf{1}, 825--835 (1933).

\bibitem{Kolos1968}
W.~Ko\l{}os and L.~Wolniewicz,
``Improved theoretical ground-state energy of the hydrogen
molecule,''
J. Chem. Phys. \textbf{49}, 404--410 (1968).

\bibitem{Kato1957}
T.~Kato,
``On the eigenfunctions of many-particle systems in quantum
mechanics,''
Commun. Pure Appl. Math. \textbf{10}, 151--177 (1957).

\bibitem{Fock1954}
V.~Fock,
``On the Schr\"odinger equation of the helium atom,''
Izv. Akad. Nauk SSSR Ser. Fiz. \textbf{18}, 161--172 (1954).

\bibitem{Morgan1986}
J.~D. Morgan~III,
``Convergence properties of Fock's expansion for S-state
eigenfunctions of the helium atom,''
Theor. Chim. Acta \textbf{69}, 181--223 (1986).

\bibitem{Numba2015}
S.~K. Lam, A.~Pitrou, and S.~Seibert,
``Numba: A LLVM-based Python JIT compiler,''
in \textit{Proceedings of the Second Workshop on the LLVM
Compiler Infrastructure in HPC} (ACM, New York, 2015),
pp.\ 1--6.

\bibitem{QUADPACK}
R.~Piessens, E.~de Doncker-Kapenga, C.~W. \"Uberhuber,
and D.~K. Kahaner,
\textit{QUADPACK: A Subroutine Package for Automatic Integration}
(Springer-Verlag, Berlin, 1983).

\end{thebibliography}

\end{document}
